\section{The Hilbert-target obstruction in the exact printed central norm}
\label{sec:printed-hilbert-cutoff-obstruction}

\paragraph{Status and scope.}
The complete post-writing author reading and two independent
whole-proof readings have passed.
This section uses the literal printed realization at twist \(A=1/2\).
The output is only assumed to belong to its actual closed theta ideal.
No finite-anticipation condition or smooth output is assumed.
The proof explicitly uses an established subset of critical-line zeros;
it does not assume \(\mathsf A\), \(M\), or the Riemann hypothesis.

Let \(\mathscr I\) denote the original rational centered smooth theta
ideal, and put
\[
 \begin{gathered}
 p_s(x)=(1+x^2)^{s/2},\qquad
 H_s^{\rm pr}=L^2(\mathbb R,p_s(x)\,dx),\qquad s>1,\\
 J_s^{\rm pr}=\overline{\mathscr I}^{\,H_s^{\rm pr}},\qquad
 \tau_yf(x)=f(x-y),\qquad \mathcal D=C_c^\infty(\mathbb R).
 \end{gathered}
 \tag{MA-PI.180}
\]
The space \(\mathcal D\) carries its usual strict LF topology.
The rational moment subspace is
\[
 \mathcal D_0=\left\{v\in\mathcal D:
                   \int_{\mathbb R}e^{-x/2}v(x)\,dx=0\right\},
 \tag{MA-PI.181}
\]
with its inherited topology. The ideal in (MA-PI.180) is the
closure of the actual theta image, not a separately postulated
spectral model.

\paragraph{The actual completed value functionals.}
For every \(\gamma\in\mathbb R\), Cauchy--Schwarz gives
\[
 \begin{aligned}
 \|h\|_{L^1}
   &\le c_s\|h\|_{H_s^{\rm pr}},&
 c_s^2&=\int_{\mathbb R}(1+x^2)^{-s/2}\,dx<\infty,\\
 \ell_\gamma(h)&=\int_{\mathbb R}h(x)e^{i\gamma x}\,dx,&
 |\ell_\gamma(h)|&\le c_s\|h\|_{H_s^{\rm pr}}.
 \end{aligned}
 \tag{MA-PI.182}
\]
If \(\rho=1/2+i\gamma\) is a nontrivial zeta zero, the original
smooth theta Mellin identity, in the normalization of
(MA-PI.151) and MT-M24, gives \(\ell_\gamma(\mathscr I)=0\).
The continuous extension (MA-PI.182) therefore annihilates
\(J_s^{\rm pr}\). This argument uses the actual inherited norm
and its actual closed ideal. It needs only a value constraint;
no derivative evaluation or unproved zero multiplicity bound
is used.

\paragraph{MA-PI-008 (no nonzero Hilbert-valued transport in the printed central norm).}
Suppose that a continuous linear map
\(T:\mathcal D\to J_s^{\rm pr}\) commutes with every translation.
Assume there is \(R\ge0\) such that, whenever the compact
supports of \(f,g\in\mathcal D\) are ordered and separated by
more than \(R\), one has
\[
       \langle T\tau_{-L}f,T\tau_{-L}g\rangle_{H_s^{\rm pr}}=0
\]
for all sufficiently large \(L\), with the threshold allowed
to depend on \(f,g\). Then \(T=0\).
The same conclusion holds for a continuous equivariant map
\(T:\mathcal D_0\to J_s^{\rm pr}\) with this property on its
domain. In particular, the exact physical-frequency plateau
and orthogonal global PW identities below cannot hold for a
nonzero such map. This applies to the printed \(A=1/2\),
\(s\ge3\) regime, without a finite-anticipation hypothesis.

\begin{proof}
\emph{1. The translated Gram form in this precise metric.}
For \(h,k\in H_s^{\rm pr}\), change variables in the inner
product to obtain
\[
 \lim_{L\to\infty}L^{-s}
       \langle\tau_{-L}h,\tau_{-L}k\rangle_{H_s^{\rm pr}}
                    =\int_{\mathbb R}h(x)\overline{k(x)}\,dx.
 \tag{MA-PI.183}
\]
Indeed \(L^{-s}p_s(x-L)\to1\), and, for \(L\ge1\),
\(L^{-s}p_s(x-L)\le C_s p_s(x)\).
Cauchy--Schwarz bounds the product by an integrable weighted
majorant. Equivariance and the assumed eventual orthogonality
therefore imply
\(\langle Tf,Tg\rangle_{L^2}=0\) for every such separated pair.
Notice that no exponential weight, equivalent norm, or limit
of different orthogonal projectors has entered this step.

\emph{2. The compact Gram distribution.}
Regard \(T\) now as a continuous map \(\mathcal D\to L^2\).
It has the translation-equivariant continuous extension
\(\widetilde T:\mathcal S(\mathbb R)\to L^2\) constructed in
MA-PI-005: use a unit-interval partition, translate its pieces
to one fixed compact interval, apply the finite derivative
bound furnished by continuity there, and sum the bounds
against Schwartz decay. The series converges absolutely in
\(L^2\); density gives equivariance. This extension is used
only in \(L^2\). Its range is not asserted to lie in
\(H_s^{\rm pr}\) or in \(J_s^{\rm pr}\).

With \(\widehat f(t)=\int f(x)e^{itx}\,dx\), the convolution
identity for \(\widetilde T\) and a Gaussian with nonvanishing
Fourier transform give a locally square-integrable multiplier
\(M\) such that
\[
 \begin{aligned}
 \widehat{\widetilde T f}(t)&=M(t)\widehat f(t),\\
 \langle Tf,Tg\rangle_{L^2}
    &=\frac1{2\pi}\int_{\mathbb R}|M(t)|^2
                        \widehat f(t)\overline{\widehat g(t)}\,dt.
 \end{aligned}
 \tag{MA-PI.184}
\]
For completeness, the convolution identity is
\(\widetilde T(f*g)=(\widetilde T f)*g=f*(\widetilde T g)\).
Its Fourier transform with one Gaussian factor proves the
multiplier identity by division by that nonzero transform.
Continuity on Schwartz space applied to modulated Gaussians
bounds the \(L^2\) mass of \(M\) in every unit frequency
interval by one fixed polynomial in its center.
Thus \(|M|^2\) is tempered.

Let \(k\) be its inverse Fourier distribution. The separated
pairings in (MA-PI.184), with one factor an approximate
identity at zero, test \(k\) against arbitrary compact tests
in either component outside \([-R,R]\). They give
\[
 \begin{gathered}
 \operatorname{supp}k\subset[-R,R],\qquad
 E=\widehat k\ \hbox{entire},\\
 |E(z)|\le C(1+|z|)^N e^{R_1|\operatorname{Im}z|}
          \quad(z\in\mathbb C)
 \end{gathered}
 \tag{MA-PI.185}
\]
for fixed \(C,N,R_1\). The last bound follows by evaluating
the compact distribution of finite order on a fixed compact
cutoff times \(e^{izx}\). Its real-axis restriction agrees
with \(|M|^2\) as a distribution. These are the actual
Gram-distribution steps of (MA-PI.156)--(MA-PI.158);
they require only the \(L^2\) continuity just established.

\emph{3. The multiplier is continuous on the real axis.}
This is the point at which the precise weight \(s>1\) is used.
For any \(t_0\in\mathbb R\), choose \(g\in\mathcal D\) with
\(\widehat g(t_0)\ne0\). The actual output \(Tg\) belongs to
\(H_s^{\rm pr}\subset L^1\), so \(\widehat{Tg}\) is continuous.
On a neighborhood where \(\widehat g\ne0\), (MA-PI.184) gives
the continuous representative
\[
 \begin{aligned}
        M(t)&=\frac{\widehat{Tg}(t)}{\widehat g(t)}
                    &&\text{where }\widehat g(t)\ne0,\\
        E(t)&=|M(t)|^2 &&(t\in\mathbb R).
 \end{aligned}
 \tag{MA-PI.186}
\]
The local representatives agree almost everywhere, hence
everywhere on their overlaps by continuity. They therefore
patch on the whole real line. The last equality follows
because both real-axis functions are continuous and their
distributions already agree. No boundary holomorphic
continuation or smoothness of the output is being assumed.

For each simple critical-line zeta zero \(1/2+i\gamma\),
choose such a compact \(g\) at \(t_0=\gamma\).
By (MA-PI.182), the actual range condition gives
\[
     \widehat{Tg}(\gamma)=\ell_\gamma(Tg)=0,
                     \qquad M(\gamma)=0,\qquad E(\gamma)=0.
 \tag{MA-PI.187}
\]
Conrey (1989), Theorem 1, p.~4, and the counting conventions
(8), (11), (12), p.~3, imply at least \(cT\log T\) distinct
such positive ordinates below \(T\), for all sufficiently
large \(T\) and some \(c>0\).
If \(E\ne0\), (MA-PI.185) and Jensen's formula, centered at
any point where \(E\ne0\), bound its number of zeros in
\(|z|\le T\) by \(O(T)\). This is a contradiction.
Thus \(E=0\), whence \(M=0\), and (MA-PI.184) gives \(T=0\).
Only compact inputs were used for the range condition in
this argument; no range property of \(\widetilde T\) was needed.

\emph{4. The exact rational moment domain.}
For a map defined on \(\mathcal D_0\), precompose with
\(D-1/2:\mathcal D\to\mathcal D_0\).
It is continuous, translation-equivariant and support
preserving. Formula (MA-PI.161) proves its surjectivity:
the primitive has compact support precisely because the
moment in (MA-PI.181) vanishes. The preceding argument
applies to this composite, so its vanishing implies
the original map is zero.
\end{proof}

\paragraph{The actual printed PW corollary.}
Here the passage to the printed source is an exact unitary
identification. Let \(m=\mu_1(C_{\mathbb Q}^1)>0\), use the
positive real norm splitting, and define
\[
 \begin{aligned}
 U:H_s^{\rm pr}&\longrightarrow
           W^{s,2}(C_{\mathbb Q})(\mathbf1),\\
 (Uh)(\nu(e^x)k)&=m^{-1/2}h(x),\\
 P_Y&=U^{-1}Q^{\rm src}_{\{\infty\},Y,1/2,s}U.
 \end{aligned}
 \tag{MA-PI.188}
\]
Masullo Definition 9.1.2.1, (9.1.2.3), p.~399, gives
exactly the norm \(p_s(x)\,dx\) after compact-fiber integration.
The centered theta image satisfies
\(|g|^{1/2}\Phi_{\mathbb Q}=\sqrt m\,U Y\).
Therefore \(U J_s^{\rm pr}\) is the actual printed closed
principal theta ideal, and (MA-PI.188) is its actual
orthogonal PW projection, as defined by Definition
11.2.1.3, (11.2.1.7--8), pp.~488--489.
This is the specialization \(b=A-1/2=0\) of the unitary
calculation (MA-PI.32)--(MA-PI.36), not a comparison of
equivalent Hilbert norms.

Use the genuine even Fourier inputs and their exact coordinate
\(v(x)=e^{-x/2}u(e^{-x})\) from (MA-PI.162).
The two original polar conditions are exactly \(u(0)=0\)
and \(v\in\mathcal D_0\). Their source membership and
dilation normalization were checked there directly from
Definition 7.3.2.3, (7.3.2.12--17), pp.~354--355.
The physical Fourier ball on p.~516, Step II, has the
eventual plateaus \(1\) for \(r\le Y/2\) and \(0\) for
\(r\ge2Y\). Thus any fixed continuous equivariant
Hilbert-valued comparison with exact plateau transport must
satisfy, for all sufficiently large \(Y\),
\[
 \begin{array}{ll}
 P_YTv=Tv,&\operatorname{supp}v\subset[c+\log2,\infty),\\
 P_YTv=0,&\operatorname{supp}v\subset(-\infty,c-\log2],
 \end{array}
                       \qquad c=-\log Y.
 \tag{MA-PI.189}
\]
For two ordered supports separated by more than \(2\log2\),
choose \(c\) between the corresponding plateaus. One image
is fixed and the other killed by the same orthogonal
projection, so they are orthogonal. Translate both inputs
far left to make the required \(Y\) arbitrarily large.
This proves the hypothesis of MA-PI-008 with \(R=2\log2\),
and hence \(T=0\). Equivalently one may start with a map
into the printed space and conjugate it by \(U^{-1}\).

The argument uses only plateau inputs, where the physical
cutoff fixes or kills them and preserves both polar conditions.
It does not claim that an arbitrary transition profile
preserves \(\mathcal D_0\). It also does not identify these
global \(S_0\) PW inputs with the coordinate-flat local
\(S_{00}\) source on p.~515.

\paragraph{Arithmetic and analytic boundaries.}
The hypothesis \(s>1\) is used exactly in (MA-PI.182).
The theorem covers the source's \(s\ge3\) regime, but makes
no assertion for \(s\le1\). The critical-line count is an
explicit known partial zero-location input, in the precise
published form already recorded in MA-MO-023. It supplies
neither a location nor a simplicity statement for the
remaining zeros. This is a cited classical result, not
a fresh proof of Conrey's mollifier theorem.

The result rules out the specified fixed equivariant maps
into the actual printed closed ideal with all-level plateau
identities. It does not assert that the empty comparison
section on p.~514 supplied such a map, that every possible
trace comparison must have one, or that a Gram-corrected,
level-dependent, or polar-augmented construction is covered.
The arithmetic FT002 identity remains a separate question.

In the inherited cote 1 placement, BR-TR-01/02,
scans079--080 and111--112, the critical value gives an actual
continuous length-one quotient of
\(H_s^{\rm pr}/J_s^{\rm pr}\).
Its continuity is the \(j=0\) case of \(2j<s-1\);
no unbounded jet is inserted. The obstruction precedes a
nuclear trace or determinant and cannot be repaired merely
by a formal finite-length or weight assignment.
